\section{Results and Verification}

\subsection{Length tuning and resonance control}
The \SI{150}{mm} analytical starting length resonated below \SI{1}{GHz}; the plotted
minimum is approximately \SI{0.83}{GHz} with $S_{11}\approx-\SI{18}{dB}$.
Shortening the model to \SI{125}{mm} moved the plotted minimum above the target to
approximately \SI{1.08}{GHz} with $S_{11}\approx-\SI{40}{dB}$. These two results
establish the expected inverse relationship between physical length and resonant
frequency.

\begin{figure}[H]
\centering
\begin{subfigure}[t]{0.48\textwidth}
\centering
\maybeincludegraphics[width=\linewidth]{figures/cst_results/fig_s11_L150_initial_halfwave.png}
\caption{SIM-150, plot-derived minimum near \SI{0.83}{GHz}.}
\end{subfigure}\hfill
\begin{subfigure}[t]{0.48\textwidth}
\centering
\maybeincludegraphics[width=\linewidth]{figures/cst_results/fig_s11_L125_tuning_run.png}
\caption{SIM-125, plot-derived minimum near \SI{1.08}{GHz}.}
\end{subfigure}
\caption{Preliminary length-tuning evidence from CST. Numerical values are approximate because raw curve exports were not supplied.}
\label{fig:preliminary-s11}
\end{figure}

\subsection{Final tuned match}
For SIM-136, the CST marker gives
\begin{equation}
f_{\mathrm{res}}=\SI{0.99648}{GHz},
\qquad
S_{11,\min}=-\SI{47.01186}{dB}.
\end{equation}
The absolute frequency error is
\begin{equation}
\Delta f=|1.00000-0.99648|=\SI{0.00352}{GHz}=\SI{3.52}{MHz},
\end{equation}
and the relative error is
\begin{equation}
\frac{\Delta f}{f_0}\times100=\SI{0.352}{\percent}.
\end{equation}

\begin{figure}[H]
\centering
\maybeincludegraphics[width=0.95\linewidth]{figures/cst_results/fig_s11_L135p6_tuned_final_with_marker.png}
\caption[Final tuned reflection coefficient with exact CST marker.]{Exact CST marker for SIM-136: $S_{11}=-\SI{47.01186}{dB}$ at \SI{0.99648}{GHz}. The equivalent positive return loss is \SI{47.01186}{dB} relative to \SI{73}{\ohm}.}
\label{fig:final-s11}
\end{figure}

\subsection{Far-field behavior}
At \SI{1}{GHz}, the final 3D plot displays a broadside toroidal pattern with axial nulls
and a peak directivity of \SI{2.149}{dBi}. The associated $\phi=90^\circ$ cut displays
a main-lobe magnitude of approximately \SI{2.15}{dBi} and a $-\SI{3}{dB}$ angular width
of approximately $78.4^\circ$.

\begin{figure}[H]
\centering
\begin{subfigure}[t]{0.58\textwidth}
\centering
\maybeincludegraphics[width=\linewidth]{figures/cst_results/fig_3d_radiation_L135p6_tuned_final.png}
\caption{3D directivity at \SI{1}{GHz}.}
\end{subfigure}\hfill
\begin{subfigure}[t]{0.39\textwidth}
\centering
\maybeincludegraphics[width=\linewidth]{figures/cst_results/fig_2d_radiation_L135p6_tuned_final.png}
\caption{$\phi=90^\circ$ principal-plane cut.}
\end{subfigure}
\caption{Final CST far-field evidence for SIM-136.}
\label{fig:final-patterns}
\end{figure}

\subsection{Analytical verification}
\begin{figure}[H]
\centering
\begin{subfigure}[t]{0.46\textwidth}
\centering
\maybeincludegraphics[width=\linewidth]{figures/analytical/analytical_half_wave_dipole_polar.png}
\caption{Normalized polar field pattern.}
\end{subfigure}\hfill
\begin{subfigure}[t]{0.50\textwidth}
\centering
\maybeincludegraphics[width=\linewidth]{figures/analytical/analytical_half_wave_dipole_pattern_db.png}
\caption{Normalized pattern in decibels and HPBW.}
\end{subfigure}
\caption{Independent analytical Python evaluation of the canonical thin half-wave dipole pattern.}
\label{fig:analytical-patterns}
\end{figure}

\begin{table}[H]
\centering
\caption{Comparison of analytical and final simulated radiation evidence.}
\label{tab:analytical-simulation-comparison}
\begin{tabularx}{\textwidth}{>{\bfseries}p{0.27\textwidth}XXX}
\toprule
Metric & Analytical & Final CST display & Interpretation \\
\midrule
Reference frequency & \SI{1.000}{GHz} & Pattern at \SI{1.000}{GHz} & Same comparison frequency \\
Peak directivity & \SI{2.15}{dBi} & \SI{2.149}{dBi} & Agreement within displayed precision \\
HPBW & $78.08^\circ$ & approximately $78.4^\circ$ & Difference about $0.32^\circ$ \\
Pattern topology & Broadside, axial nulls & Broadside, axial nulls & Canonical dipole behavior \\
\bottomrule
\end{tabularx}
\end{table}

\subsection{Verification and validation status}
The analytical script independently verifies the canonical pattern metrics. The CST
figures verify the qualitative tuning trend and final simulated response. Experimental
validation is absent; therefore, the package demonstrates analytical-to-simulation
agreement rather than hardware performance.
